\subsection{\texorpdfstring{关于空间 \(L^{p}(X,\mu)\) 的补充}{关于空间 L\^{}\{p\}(X,\textbackslash mu) 的补充}}

命题 1. --- 设 X 是局部紧拓扑空间。设 \(\mu\) 是 X 上的正测度，\(p \in [1, +\infty[\)。设 \((\mathrm{X}_i)_{i \in \mathrm{I}}\) 是 X 中两两不相交的局部紧子集的局部可数族（INT, IV, p. 190, § 5, no. 9，定义 7），使得诸 \(\mathrm{X}_i\) 之并的余集是局部 \(\mu\)-零测的。设 \(\varphi_i\) 是 \(\mathrm{X}_i\) 的特征函数，\(\mu_i\) 是 \(\mu\) 在 \(\mathrm{X}_i\) 上诱导的测度（INT, IV, p. 186, §5, n° 7，定义 4）。

\begin{enumerate}
\def\labelenumi{\alph{enumi})}
\item
  映射 \(p_{i}\) ：\(f \mapsto f\varphi_{i}\) 是 \(\mathcal{L}^{p}(\mathrm{X}, \mu)\) 的投影子，且过渡到商，定义了投影子 \(\widetilde{p}_{i}\) （\(\mathrm{L}^{p}(\mathrm{X}, \mu)\) 的投影子）。当 p = 2 时，后者是 \(\mathrm{L}^{2}(\mathrm{X}, \mu)\) 的正交投影子；
\item
  映射 \(f \mapsto f|X_i\) 诱导出从 \(p_i\) 的像到空间 \(\mathcal{L}^p(X_i, \mu_i)\) 上的等距同构，且过渡到商，定义了从 \(\widetilde{p}_i\) 的像到 \(\mathrm{L}^p(X_i, \mu_i)\) 上的等距同构，借助它这两个空间互相等同；
\item
  诸空间 \(\mathrm{L}^{p}(\mathrm{X}_{i},\mu_{i})\) 之和在空间 \(\mathrm{L}^{p}(\mathrm{X},\mu)\) 中是完全的。当 p = 2 时，空间 \(\mathrm{L}^{2}(\mathrm{X},\mu)\) 是诸空间 \(\mathrm{L}^{2}(\mathrm{X}_{i},\mu_{i})\) 的希尔伯特和。
\end{enumerate}

断言 \(a\)）是初等的。断言 \(b\)）由 INT, V, p. 84, § 7, n° 1 的评注得出。

设 \(q \in ]1, +\infty]\) 是 \(p\) 的共轭指数。我们把 \(L^p(X, \mu)\) 的对偶与 \(L^q(X, \mu)\) 等同起来（INT, V, p. 61, § 5, n° 8，定理 4）。设 \(f\) 是 \(\mathcal{L}^q(X, \mu)\) 中的函数，其类 \(\widetilde{f}\) 在 \(L^q(X, \mu)\) 中满足 \(\langle \widetilde{f}, \widetilde{p}_i(\varphi) \rangle = 0\) （对所有 \(i \in I\) 与每个 \(\varphi \in L^p(X, \mu)\)）。由于 \(p_i\) 的像包含 \(\mathcal{K}(X_i)\) ，可得测度 \((f|X_i) \cdot \mu_i\) 在 \(X_i\) 上为零，即 \(f\) 在 \(X_i\) 上的限制关于 \(\mu_i\) 在 \(X_i\) 上局部零测（INT, V, p. 46, § 5, n° 3，推论 2）。由于诸 \(X_i\) 之并在 X 中的余集是局部 \(\mu\)-零测的，可以断定函数 \(f\) 在 X 上局部 \(\mu\)-零测（INT, IV, p. 190, § 5, n° 9 及 p. 172, n° 2，命题 5）。于是 \(f\) 的类在 \(L^q(X, \mu)\) 中为零（当 \(p \neq 1\) 时用到 INT, V, p. 8, § 1, n° 3 的引理 1 及命题 9 的推论）。由哈恩--巴拿赫定理（EVT, II, p. 46，推论 1），断言 \(c\)）的第一部分得证。

若 \(p = 2\) ，\(p_i\) 的像与 \(p_j\) 的像正交（对所有 \(i \neq j\) ），因为此时 \(X_i\) 与 \(X_j\) 不相交，由此得到最后一个断言。

命题 2. --- 设 X 是在无穷远处可数的局部紧空间。设 \(\mu\) 是 X 上的正测度。对每个 \(p \in [1, +\infty[\)，空间 \(\mathrm{L}^{p}(\mathrm{X},\mu)\) 是可数型的。

设 \((\mathrm{U}_{n})_{n\in\mathbf{N}}\) 是 X 中相对紧开集的序列，其并等于 X 且满足 \(\overline{U}_{n}\subset U_{n+1}\) 对所有 \(n\in\mathbf{N}\) 成立（TG, I, p. 68，命题 15）。对每个 \(n\in\mathbf{N}\)，空间 \(\mathcal{K}(\mathrm{X},\overline{\mathrm{U}}_{n})\) 等同于巴拿赫空间 \(\mathcal{C}(\overline{\mathrm{U}}_{n})\) 的闭子空间（INT, III, p. 40, § 1, n° 1）；由于后一空间是可数型的（TG, X, p. 25，推论），\(\mathcal{K}(\mathrm{X},\overline{\mathrm{U}}_{n})\) 亦然（TG, IX, p. 19, cor., (i)）。设 \(F_{n}\) 是 \(\mathcal{K}(\mathrm{X},\overline{\mathrm{U}}_{n})\) 的可数稠密子集。

设 \(f \in \mathcal{L}^p(\mathrm{X},\mu)\) 且 \(\varepsilon > 0\) 。存在整数 \(n \in \mathbf{N}\) 使得 \(\int_{\mathrm{X}-\mathrm{U}_n}|f|^p < \varepsilon / 2\) ，且存在 \(g \in \mathcal{F}_n\) 使得 \(\int_{\overline{\mathrm{U}}_n}|f - g|^p < \varepsilon / 2\) 。因此，\(\mathrm{L}^p(\mathrm{X},\mu)\) 中诸集合 \(\mathcal{F}_n\) 的元素之类所成的并在 \(\mathrm{L}^p(\mathrm{X},\mu)\) 中稠密，证明完成（TG, IX, p. 18，命题 12）。

